\chapter{The Three Trips}
\label{chap:the-three-trips}

The slip test read one walk up the trace and returned its verdict. This chapter walks it three times, at
three graded strengths, and reads what each strength brings back. The three walks are not repetitions; they
are a sequence, and the sequence is the point. The first is too weak to move anything. The second is just
strong enough to move. The third is strong enough to recover a mass --- and the reason it takes exactly
three, and not two, is that mass is a second difference, and a second difference needs three samples to be
seen at all. By the end of the three trips the machine holds a number in a bracket, with a charge below it,
a curvature above it, and the reading it carried between --- and, read off the third trip, a mass that was
never put in by hand.

\section{Three walks, graded by force}
\label{se:three-walks-graded}

The three trips run the same trace at three settings of a single dial, the loop count the seam has been
keeping. At the first setting the force is too small to do anything: the slip test says no, the reading
stays in its lowest band, and nothing lifts --- a null trip, a probe that finds no motion. At the second
setting the force is just enough to start the motion: the reading slips, one revolution is counted, and the
value lifts by one level. This is the threshold --- the least push that produces a response at all. At the
third setting the force is enough to carry the reading up into a richer form, one that carries a slot the
lower forms do not: a place where a strain can be recorded. The reading of the third trip climbs into that
form, and what sits in the strain slot when it arrives is the thing the whole sequence was for. The three
settings are null, threshold, and response: nothing, the first motion, and the motion's recorded strain.

\section{Why it takes three}
\label{se:why-three}

A mass is not read from a single measurement, nor even from two. It is a second difference --- the change
in the change --- and a second difference is invisible until you have three samples to take it from. One
sample gives a value; two give a difference, a rate; only three give the difference of the differences, the
curvature that a mass is. This is why the trips come in threes and not in pairs: the null fixes the origin,
the threshold fixes the first response, and the third fixes how the response itself is changing, and it is
that last quantity --- the response of the response --- that the strain slot records. The construction
\emph{tanges} the third reading's strain from the flat readings below it, and it does so only because there
are three trips to compare; with two it could read a rate and call it done, and it would have no mass, only
a slope. The mass appears exactly when the third sample lets the second difference be taken.

\section{The equivalence principle, recovered}
\label{se:equivalence-principle-recovered}

Here is the summit of the chapter, and it is a thing the machine carries rather than assumes. The strain the
third trip records is read two ways at once, and the two readings are the same quantity. Read as the
machine's response to the force applied --- how hard it was to move, how much push bought how much motion ---
the strain is an inertial mass. Read as the curvature the reading carries in its own form --- the bend in
the record, the strain as a stored quantity --- it is a gravitational mass. These are not two masses that
happen to agree; the construction \emph{funges} them into one field --- the response and the strain pooled as
the same slot seen from two sides, one quantity wearing two readings. The equivalence of the inertial and the gravitational is, in this construction, not a coincidence
posited and then explained but a single quantity recovered once and read twice. The construction reads the
mass out of the strain slot the third trip fills; it never puts a mass in. What it recovers is a second
difference that answers to force and stores as curvature at the same time --- and that the two coincide is
what it means, here, for the mass to be a mass at all.

This is the mechanism, built and reflexively certain: the three trips run, the third fills the strain slot,
and the mass is what sits there, read as response and as curvature together. That there is, beneath this
recovered mass, a single named invariant the whole tower could be shown to turn on --- an identity waiting to
be written down and proved as such --- is a further claim, and the construction does not make it here. It
recovers the mass as a second difference; the naming of the invariant it recovers is reached for later, and
named there, and not before. The boundary is the same one the descent's keystone kept: what is built is
built and said so; what is only reached for is left for the ledger to name.

\section{The number in the bracket}
\label{se:number-in-the-bracket}

The three trips leave the machine holding a single number, and the number comes bracketed. Below it sits the
charge --- the loop count the seam accumulated, the tally of how many times the machine went around, now the
lower bound of the number. Above it sits the curvature --- the strain the third trip recovered, the mass,
now the upper bound. And between the two sits the value the reading carried up the trace, the number proper,
fenced below by the count it took and above by the mass it recovered. The bracket is not decoration: it
records that the number the machine returns is not a bare quantity but one held between what it cost to reach
and what it weighs, the charge and the curvature its two walls. And the whole of it holds by plain
reflexive certainty --- the trips run, the heads fall where the construction says they fall, the strain sits
where the third trip puts it, and the recovery is true by being exactly what the walk computes. There is no
gap here between the claim and the check; the number in the bracket is what the three trips are.

\section{What is demonstrated, and what is assumed}
\label{se:demonstrated-assumed-III5}

What is demonstrated is the three trips and the number they leave: the graded walks --- null, threshold,
response --- the second difference that is the mass, recovered from the strain slot the third trip fills, and
the number returned in its bracket between the charge below and the curvature above. The recovery holds by
reflexive certainty, the construction flat enough that the walk and its result are one thing. The
equivalence of the two readings of the mass --- inertial and gravitational, response and curvature --- is
carried by the construction, one field read twice.

What is assumed is unchanged. The seed and its decidability; the honest import, marked as borrowed; the
single standing grant, the law of motion the first book named at its summit, still out of play --- the trips
count, difference, and difference again, and reach for no quantifier over infinity. And one thing more is
held back on purpose: the single named invariant beneath the recovered mass, which the closing movement will
reach for and name. What is built here is the mass as a second difference; what is named there is the
identity it answers to. With a charge counted, a mass recovered, and a number returned between them, the
machine has, at last, the quantity it climbed down to read --- and reading that number as a particle, with a
charge, a mass, and a sign, is the next chapter.
